\section{결론}

이 논문은 같은 경주에 열린 여러 패리뮤추얼 시장을 하나의 순위 상태공간 위에
놓는다. 삼쌍승식은 상위 3마리의 순서를 구분하는 가장 세분된 청구권이며, 다른
주요 승식은 그 상태들을 합한 사건에 지급한다. 따라서 승식별 가격의 관계를
단순한 상관관계가 아니라 명시적인 가산 제약으로 검정할 수 있다.

이 설계가 제공하는 핵심 구분은 내부 일관성과 외부 정확도의 분리다. 삼쌍승식
가격이 단순 승식 가격을 잘 재구성한다면 여러 풀이 공통된 순위 구조를 반영한다는
증거가 된다. 그러나 그것만으로 그 구조가 객관적 확률과 일치한다고 말할 수는
없다. 실현 착순에 대한 예측력과 전통적인 인기마--비인기마 편향은 별도 검정으로
확인해야 한다. 이 구분을 유지하면 경마시장을 단순한 이상현상의 사례가 아니라,
서로 다른 차원의 시장이 정보를 집계하고 전달하는 방식을 관찰할 수 있는
실험장으로 사용할 수 있다.
